\documentclass[12pt,oneside]{book}

% ============================================================
% METADATOS
% ============================================================
\newcommand{\universidad}{Universidad de Buenos Aires (UBA)}
\newcommand{\facultad}{Facultad de Derecho}
\newcommand{\carrerauniv}{Carrera de Abogacía}
\newcommand{\materia}{Derecho Procesal Civil y Comercial}
\newcommand{\titulo}{El Proceso de Desalojo: Diligencias Preliminares y Medidas Preparatorias}
\newcommand{\autor}{Isaac Edgar Camacho Ocampo}
\newcommand{\anio}{2026}

% ============================================================
% PAQUETES
% ============================================================
\usepackage[utf8]{inputenc}
\usepackage[T1]{fontenc}
\usepackage[spanish,es-nodecimaldot]{babel}

\usepackage[
    top=2.5cm,
    bottom=2.5cm,
    left=3cm,
    right=2.5cm,
    headheight=15pt
]{geometry}

\usepackage{amsmath}
\usepackage{amssymb}
\usepackage{mathtools}

\usepackage{graphicx}
\usepackage{float}
\usepackage{caption}
\usepackage{subcaption}

\usepackage{booktabs}
\usepackage{array}
\usepackage{longtable}
\usepackage{tabularx}

\usepackage{enumitem}

\usepackage{xcolor}
\usepackage[most]{tcolorbox}

\usepackage{fancyhdr}
\usepackage{ragged2e}

\usepackage[
    colorlinks=true,
    linkcolor=blue!50!black,
    citecolor=green!40!black,
    urlcolor=blue!60!black,
    pdftitle={\titulo},
    pdfauthor={\autor},
    pdfsubject={Apuntes universitarios},
    bookmarks=true
]{hyperref}

% ============================================================
% CONFIGURACIÓN
% ============================================================
\graphicspath{
    {./img/}
    {./graficos/}
    {./figuras/}
}

\setcounter{tocdepth}{2}

\setlength{\parindent}{0pt}
\setlength{\parskip}{0.5em}

\setlist[itemize]{
    topsep=0.3em,
    itemsep=0.2em
}

\setlist[enumerate]{
    topsep=0.3em,
    itemsep=0.2em
}

\clubpenalty=10000
\widowpenalty=10000

\pagestyle{fancy}
\fancyhf{}
\fancyhead[L]{\nouppercase{\leftmark}}
\fancyhead[R]{\materia}
\fancyfoot[C]{\thepage}
\renewcommand{\headrulewidth}{0.4pt}
\renewcommand{\footrulewidth}{0pt}

\fancypagestyle{plain}{
    \fancyhf{}
    \fancyfoot[C]{\thepage}
    \renewcommand{\headrulewidth}{0pt}
}

% ============================================================
% COMANDOS
% ============================================================
\newcommand{\abs}[1]{\left|#1\right|}
\newcommand{\norm}[1]{\left\lVert#1\right\rVert}
\newcommand{\vect}[1]{\mathbf{#1}}
\newcommand{\deriv}[2]{\frac{d#1}{d#2}}
\newcommand{\pderiv}[2]{\frac{\partial#1}{\partial#2}}

% ============================================================
% ENTORNOS / CAJAS
% ============================================================
\newtcolorbox{definition}[1]{
    enhanced,
    breakable,
    title={Definición: #1},
    fonttitle=\bfseries,
    colback=blue!3,
    colframe=blue!50!black,
    boxrule=0.8pt,
    arc=2mm
}

\newtcolorbox{important}{
    enhanced,
    breakable,
    title={Importante},
    fonttitle=\bfseries,
    colback=yellow!8,
    colframe=orange!70!black,
    boxrule=0.8pt,
    arc=2mm
}

\newtcolorbox{example}[1]{
    enhanced,
    breakable,
    title={Ejemplo: #1},
    fonttitle=\bfseries,
    colback=green!4,
    colframe=green!40!black,
    boxrule=0.8pt,
    arc=2mm
}

\newtcolorbox{formula}[1]{
    enhanced,
    breakable,
    title={Fórmula / Regla: #1},
    fonttitle=\bfseries,
    colback=purple!4,
    colframe=purple!50!black,
    boxrule=0.8pt,
    arc=2mm
}

\newtcolorbox{exercise}[1]{
    enhanced,
    breakable,
    title={Ejercicio: #1},
    fonttitle=\bfseries,
    colback=gray!5,
    colframe=gray!60!black,
    boxrule=0.8pt,
    arc=2mm
}

\newtcolorbox{note}{
    enhanced,
    breakable,
    title={Nota},
    fonttitle=\bfseries,
    colback=white,
    colframe=gray!50,
    boxrule=0.6pt,
    arc=2mm
}

\begin{document}

% ============================================================
% PORTADA
% ============================================================
\begin{titlepage}
\thispagestyle{empty}

\begin{center}

\vspace*{1cm}

{\large \textsc{\universidad}}

\vspace{0.2cm}

{\large \textsc{\facultad}}

\vspace{0.2cm}

{\large \carrerauniv}

\vspace{3cm}

{\Huge \textbf{\materia}}

\vspace{1cm}

{\LARGE \titulo}

\vspace{0.8cm}

\textit{Comprender el proceso antes de transitarlo es la mejor forma de no perderse en el camino.}

\vspace{1.2cm}

\rule{0.70\textwidth}{0.5pt}

\vspace{0.5cm}

{\large Apuntes de estudio}

\vspace{0.5cm}

\rule{0.70\textwidth}{0.5pt}

\vfill

{\large \autor}

\vspace{0.3cm}

{\large \anio}

\end{center}

\vfill

\begin{flushleft}
\textit{Este es un modesto aporte a los estudiantes de ``Derecho Procesal Civil y Comercial'' no pretende sustituir el material oficial de la materia.}
\end{flushleft}

\end{titlepage}

% ============================================================
% ÍNDICE
% ============================================================
\tableofcontents

% ============================================================
% CONTENIDO
% ============================================================

\chapter{Marco introductorio y conceptual}

Este capítulo reúne el marco de presentación de la clase: el recurso pedagógico utilizado para introducir el tema, el concepto general de diligencia preliminar y su ubicación normativa, y un caso práctico inicial que ilustra su utilidad concreta.

La clase desarrolla el instituto de las diligencias preliminares y medidas preparatorias reguladas en el Capítulo II del Libro Segundo del Código Procesal Civil y Comercial de la Nación (arts.\ 323 y siguientes), en su vinculación directa con el proceso de desalojo de inmuebles urbanos. Se abordan, entre otros ejes, la diferencia entre diligencia preliminar y prueba anticipada; los tipos de proceso de conocimiento aplicables al desalojo en el orden nacional y en el de la provincia de Buenos Aires; la legitimación sustancial y procesal; el fuero de atracción en los procesos universales y su límite respecto de las locaciones con destino habitacional; el régimen de la intimación de pago previa al desalojo por falta de pago; la prohibición de recusar sin causa en el juicio de desalojo; y un conjunto de herramientas prácticas y métodos alternativos de resolución de conflictos aplicados a la locación urbana.

\section{El arte como herramienta pedagógica del proceso}

Antes de ingresar al tema del día se retomó una dinámica ya utilizada en una clase anterior: la presentación de una obra de arte como disparador conceptual. Se trata de un cuadro del artista neerlandés Johannes Vermeer, fechado en 1665, conocido popularmente como ``la Mona Lisa del Norte'', aunque su título original es \textit{La joven de la perla}, una pintura que se exhibe en un museo de La Haya, en los Países Bajos.

Se invitó a los estudiantes a describir lo que observaban en la imagen. Entre las respuestas surgieron: una figura que ``nos está mirando con mucha luz'', un pañuelo en la cabeza, un fondo completamente negro y un arito o perla en la oreja de la joven retratada.

\begin{example}{La joven de la perla}
El docente explicó por qué eligió nuevamente esta imagen para introducir el tema del desalojo: el fondo negro del cuadro representa aquello que no se ve a simple vista en un proceso judicial —un trasfondo que hay que aprender a observar—, mientras que la perla o el arito remite a un elemento cuyo significado exacto ni siquiera puede confirmarse con el propio artista, ya fallecido, y que por eso mismo habilita el debate. El docente señaló que el fondo negro se utiliza también en psicología y en psiquiatría para referirse a aquello que el sujeto oculta, y que, trasladado al proceso judicial, es tarea del operador jurídico visualizarlo y colaborar en su develamiento.
\end{example}

La conclusión pedagógica de esta introducción fue explícita: así como detrás de la imagen de la joven del cuadro hay una historia que la pintura no revela por completo, detrás de cada proceso de desalojo hay un sujeto, una historia y circunstancias que exceden lo que el expediente muestra a primera vista, y que el abogado debe aprender a indagar.

\section{Diligencias preliminares y medidas preparatorias: marco general}

El docente aclaró, como primera precisión conceptual, que las expresiones ``diligencias preliminares'' y ``medidas preparatorias'' no son sinónimos y que conviene no confundirlas.

Se ubicó el instituto normativamente: las diligencias preliminares se encuentran reguladas en el Código Procesal Civil y Comercial de la Nación, a partir del artículo 324 y siguientes, dentro del Capítulo II (``Diligencias preliminares'') del Título I del Libro Segundo, dedicado a los procesos de conocimiento.

Consultado el grupo sobre si estas diligencias se usan habitualmente en la práctica del desalojo, el docente respondió con una precisión importante:

\begin{quote}
\textit{Docente:} En la práctica se usan poco, pero serían útiles. Serían útiles para poder focalizar el objeto de nuestro proceso de desalojo.
\end{quote}

Se explicó que un mismo inmueble puede ser objeto de distintos procesos simultáneos o sucesivos —un desalojo, una acción reivindicatoria, una diligencia preliminar—, algunos de los cuales pueden quedar abandonados y otros no, y que la diligencia preliminar se inicia siempre antes del proceso principal, con la finalidad puntual de conocer con precisión el objeto que se va a demandar. Esa precisión previa evita que, al momento de darse traslado de la demanda, el demandado oponga una defensa de falta de legitimación sustancial para estar en el proceso, por haberse demandado incorrectamente.

\section{Caso práctico inicial: la sobrina heredera}

Para ilustrar la utilidad práctica de la diligencia preliminar, el docente relató —en tercera persona, sin identificar a las partes reales— un caso reciente de su propia experiencia profesional.

\begin{example}{La sobrina heredera única}
El caso involucra a una sobrina, heredera única de un tío fallecido, soltero y sin hijos, que llevaba una vida ordenada y reservada. Al fallecer, la sobrina no encontró documentación alguna que acreditara la situación de varios inmuebles que el tío aparentemente poseía: no se sabía si estaban alquilados, prestados, o en qué carácter se encontraban ocupados por terceros. Consultada sobre datos elementales —un número de teléfono del ocupante, la existencia de una caja de seguridad, algún amigo cercano al tío—, la sobrina no pudo aportar ninguno. El perfil del causante permitía descartar razonablemente que hubiera prestado el inmueble sin contrato, pero no alcanzaba para saber si existía un boleto de compraventa, una seña, un comodato, un contrato de locación o incluso una donación.
\end{example}

Frente a este vacío de información, el docente explicó el protocolo que suele aplicar antes de recurrir a la vía judicial: intentar, en primer lugar, un contacto directo y no conflictivo con el ocupante, mediante lo que denomina informalmente ``el papelito'' —una citación no formal a los fines de un acuerdo, entregada personalmente, en la que se ofrecen múltiples canales de contacto (teléfono, correo electrónico, videollamada)—, evitando así una carta documento que, según explicó, muchas veces ni siquiera es retirada del correo.

\begin{quote}
\textit{Docente:} Yo lo que intento primero es hablar con la persona. Si usted le manda una carta documento, esa persona se va a comunicar con un abogado y el abogado con usted. Todo esto que conté es previo, que no es obligatorio, es de personalidad mía que hago eso.
\end{quote}

Ante la falta de respuesta o de apertura de la puerta por parte del ocupante, el docente recurre entonces a la diligencia preliminar prevista en el Código Procesal: solicita al juzgado que ordene al oficial de justicia constituirse en el inmueble a llevar adelante un interrogatorio previamente redactado por el propio abogado. Si nadie abre la puerta, el oficial de justicia deja constancia de ello al devolver el mandamiento, y el juez puede tener por cierto el contenido de la diligencia, habilitando así el inicio del proceso de desalojo por la causal correspondiente (falta de pago, vencimiento de contrato, comodato vencido, etc.).

\begin{definition}{Informe de dominio vs.\ índice de titularidad}
El informe de dominio es aquel que se solicita cuando ya se cuenta con el número de matrícula del inmueble en cuestión. Cuando, en cambio, lo que se busca es determinar qué inmuebles figuran a nombre de una persona determinada —sin conocer previamente la matrícula—, el trámite correcto es el índice de titularidad. Este último puede solicitarse en el registro correspondiente a cada jurisdicción, cada una con su propio organismo registral y su propio costo. Solo puede ser solicitado por profesionales habilitados: abogados, escribanos o martilleros.
\end{definition}

El docente advirtió también sobre los límites prácticos de esta búsqueda: si el heredero sospecha que el causante poseía bienes en otra provincia, no existe una vía sencilla y directa para relevar todo el país; en tales casos solo cabría solicitar al juez que oficie al registro de la propiedad de esa provincia puntual, siempre que exista algún indicio razonable que justifique la medida, ya que el juez no libraría oficios de manera genérica a las veinticuatro provincias.

\chapter{Régimen procesal de las diligencias preliminares}

\section{El artículo 323 del CPCCN y la declaración jurada}

Retomando el caso de la sobrina heredera, el docente explicó que el inciso 1º del artículo 323 del Código Procesal Civil y Comercial de la Nación habilita exactamente la medida que necesitaba en ese caso: conocer, con carácter previo a la demanda, un dato relativo a la personalidad de quien va a ser demandado.

\begin{important}
\textbf{Art.\ 323, inciso 1º, Código Procesal Civil y Comercial de la Nación.} ``El proceso de conocimiento podrá prepararse pidiendo el que pretenda demandar, o quien, con fundamento prevea que será demandado: 1) Que la persona contra quien se proponga dirigir la demanda preste declaración jurada, por escrito y dentro del plazo que fije el juez, sobre algún hecho relativo a su personalidad, sin cuya comprobación no pueda entrarse en juicio.''

% TODO: contenido incompleto o ambiguo en las notas originales
Nota de cátedra: el propio artículo 323 continúa, en sus incisos siguientes, enumerando otras diligencias preliminares posibles (exhibición de la cosa mueble, exhibición de testamento, exhibición de títulos en caso de evicción, exhibición de documentación societaria, entre otras), no todas aplicables al proceso de desalojo.
\end{important}

El docente explicó el mecanismo práctico: se trata de una declaración jurada por escrito que el propio abogado redacta con antelación, en formato de preguntas y respuestas, conforme a su mejor criterio sobre lo que resulta útil para el desalojo (por ejemplo, en qué carácter ocupa el inmueble la persona notificada). Aunque el ocupante no firme esa declaración —porque no atiende al oficial de justicia o no abre la puerta—, el contenido puede ser tenido por cierto por el juzgado.

Se hizo especial hincapié en un principio procesal que atraviesa todo el instituto: el desalojo debe dirigirse siempre contra el sujeto identificado como ocupante, subinquilino u ocupante en general, ya que de ello depende la legitimación pasiva. En todos los casos debe notificarse fehacientemente —mediante cédula o acta notarial— a quien ocupa o subocupa el inmueble, de forma tal que quede satisfecho el principio de defensa en juicio y nadie pueda luego alegar la nulidad de todo lo actuado por falta de notificación.

\section{Diligencia preliminar y prueba anticipada: dos institutos distintos}

El docente marcó una distinción que consideró central y que suele confundirse: la diligencia preliminar no debe mezclarse con el instituto de la prueba anticipada.

\begin{quote}
\textit{Docente:} El instituto de la prueba anticipada es un cajón del placar que es la diligencia, las previas o preparatorias. Son cosas diferentes, no se mezclen.
\end{quote}

\begin{definition}{Diligencia preliminar vs.\ prueba anticipada}
La diligencia preliminar tiene carácter restrictivo y residual: está limitada a los supuestos taxativos o, según cierta doctrina, enunciativos del artículo 323 y siguientes, y es restrictiva porque queda sujeta a un plazo de caducidad de treinta días para iniciar el proceso principal.

La prueba anticipada, en cambio, es un sistema asegurativo probatorio: una herramienta que el legislador otorga a las partes para no perder una prueba que podría desaparecer antes de llegar a la etapa procesal correspondiente —el ejemplo típico mencionado en clase fue el de un testigo próximo a fallecer—. Se trata, en palabras del docente, de ``una mega cautelar'', y de un instituto autónomo, no subordinado a la diligencia preliminar.
\end{definition}

\section{El plazo de caducidad de treinta días}

Vinculado directamente con el carácter restrictivo señalado en la sección anterior, el docente precisó el plazo de caducidad que rige la diligencia preliminar: una vez sustanciada, el interesado cuenta con treinta días para iniciar el proceso principal. Si no lo hace, la diligencia pierde sentido y sus efectos caducan.

\begin{quote}
\textit{Docente:} Una vez que hacemos la diligencia, tengo 30 días para empezar el proceso principal. Si no, no tiene sentido esta diligencia que hice.
\end{quote}

Ante la observación de una estudiante en el sentido de que no siempre es necesario recurrir a las diligencias preliminares, el docente coincidió, pero remarcó que, cuando faltan datos imprescindibles para no ``perderse'' en el proceso, la diligencia debe usarse y no dejarse ``enterrada'':

\begin{quote}
\textit{Estudiante:} No siempre hay que hacer las diligencias preliminares.

\textit{Docente:} Claro, si no la dejás ahí enterrada. Las tenés que usar, si faltan datos, para no marearte.
\end{quote}

\section{Derecho comparado: discovery y disclosure anglosajones}

El docente introdujo una perspectiva de derecho comparado, recuperando referencias ya utilizadas en una clase anterior sobre los institutos del \textit{discovery} (sistema estadounidense) y el \textit{disclosure} (sistema del Reino Unido), en esencia el mismo tipo de trabajo preparatorio que en Argentina se canaliza a través de las diligencias preliminares, aunque con una diferencia estructural relevante.

\begin{note}
El docente citó un artículo del doctor Salgado —presentado como compañero de cátedra—, publicado en la revista Derecho Procesal, de aproximadamente cinco páginas, que desarrolla esta comparación. Señaló que conocer el derecho comparado tiene un valor que excede lo puramente académico: ``El día de mañana ustedes pueden ser legisladores o legisladoras. Entonces uno puede proponer cambiar un poco nuestro sistema antiguón y ayornarlo un poco con lo que sirve en el derecho comparado''.
\end{note}

La diferencia central es que, en la Argentina, las diligencias preliminares son netamente judiciales —tienen que tramitarse ante un juzgado—, mientras que en los sistemas norteamericano y británico existe una etapa extrajudicial de mayor amplitud. Con todo, el docente advirtió contra la tentación de trasplantar mecánicamente institutos extranjeros:

\begin{quote}
\textit{Docente:} Aplicar lo que aplican ellos con mentalidad nuestra no va a funcionar. Eso es lo que yo digo, pero para que ustedes tengan notas de color y no tengamos todo lo mismo.
\end{quote}

\chapter{Tipos de proceso, competencia y legitimación}

\section{Procesos de conocimiento aplicables al desalojo}

En el orden nacional existen dos procesos de conocimiento: el ordinario y el sumarísimo. La gran mayoría de los procesos de desalojo tramitados en el orden nacional adquiere la vía del proceso sumarísimo, siendo minoritarios los que tramitan por la vía ordinaria. En los departamentos judiciales de la provincia de Buenos Aires, en cambio, subsisten los tres procesos de conocimiento —ordinario, sumario y sumarísimo—, lo que exige al abogado litigante conocer con precisión en qué jurisdicción actúa.

\begin{table}[H]
\centering
\caption{Procesos de conocimiento aplicables al desalojo}
\begin{tabularx}{\textwidth}{>{\raggedright\arraybackslash}X >{\raggedright\arraybackslash}X >{\raggedright\arraybackslash}X}
\toprule
\textbf{Jurisdicción} & \textbf{Procesos de conocimiento disponibles} & \textbf{Vía habitual en el desalojo} \\
\midrule
Orden nacional (Capital Federal) & Ordinario y sumarísimo & Sumarísimo (mayoría de los casos) \\
Provincia de Buenos Aires & Ordinario, sumario y sumarísimo & Depende del criterio judicial en la primera resolución \\
\bottomrule
\end{tabularx}
\end{table}

Tratándose de una materia procesal —de competencia no delegada al Congreso Nacional—, cada provincia argentina cuenta con su propio código procesal local, de modo que en algunas jurisdicciones el desalojo tramita como proceso monitorio, en otras como proceso abreviado y en otras, como en la Nación y en la provincia de Buenos Aires, como proceso de conocimiento pleno. El docente mencionó, a título ilustrativo, que en la provincia de Córdoba el proceso de desalojo tramita por vía abreviada, con una impronta particular vinculada a la perspectiva de género cuando interviene una mujer en situación de vulnerabilidad.

\begin{important}
El docente advirtió expresamente sobre una fuente frecuente de error práctico: el propio Código Procesal Civil y Comercial de la Nación establece, en su texto histórico, que el proceso de desalojo se rige por las reglas del proceso sumario; sin embargo, esa remisión se encuentra derogada, ya que el proceso sumario fue eliminado como categoría autónoma del ordenamiento procesal nacional por la Ley 25.488 (año 2001). El operador jurídico debe actualizar permanentemente su lectura del Código con cada nueva legislación que modifica el texto originario.
\end{important}

\section{El criterio judicial para fijar el tipo de proceso}

Frente a la pregunta de qué criterio sigue el juez para decidir, en la provincia de Buenos Aires, si un desalojo tramitará como ordinario, sumario o sumarísimo, el docente explicó que esa decisión se adopta en la primera resolución dictada en el expediente, y que conviene identificarla con claridad desde el inicio:

\begin{quote}
\textit{Docente:} Yo en el estudio le digo: lo ponés como un cartel en la carpeta papel. ¿Qué proceso es? Porque no tiene nada que ver los plazos en el sumarísimo con un ordinario.
\end{quote}

En cuanto al criterio sustancial que orienta esa decisión, cuando se trata de una situación simple, sin mayor conflictividad, el juez suele encausar el proceso por la vía sumarísima, que cuenta con plazos más breves y menos planteos posibles. Cuando la demanda presenta una explicación más compleja, el juez tiende a encausarla por la vía ordinaria, que habilita mayor debate.

\section{Legitimación sustancial y legitimación procesal}

\begin{definition}{Legitimación sustancial y legitimación procesal}
La legitimación sustancial (equivalente, en la terminología clásica, a la capacidad de derecho) refiere a ser el titular de la relación jurídica sustancial que se pretende hacer valer en el proceso. En un contrato de locación, corresponde a quien figura como locador o locatario en el contrato, y no a un tercero ajeno a esa relación.

La legitimación procesal (equivalente a la capacidad de hecho o de ejercicio) refiere a la aptitud para estar por sí mismo en el proceso judicial. Puede ocurrir que una persona sea titular sustancial del derecho —esté correctamente legitimada en ese plano— pero carezca de capacidad de ejercicio para actuar personalmente en el proceso, en cuyo caso deberá comparecer a través de su representante legal.
\end{definition}

\begin{important}
\textbf{Art.\ 24, Código Civil y Comercial de la Nación — Personas incapaces de ejercicio.} ``Son incapaces de ejercicio: a) la persona por nacer; b) la persona que no cuenta con la edad y grado de madurez suficiente, con el alcance dispuesto en la Sección 2ª de este Capítulo; c) la persona declarada incapaz por sentencia judicial, en la extensión dispuesta en esa decisión.''

Nota de cátedra: el docente utilizó como ejemplo el caso de un locatario que, luego de haber firmado el contrato, es declarado insano por sentencia judicial: en tal supuesto, la persona conserva su legitimación sustancial (sigue siendo la titular del contrato), pero pierde la legitimación procesal, de modo que debe comparecer en el proceso su curador.
\end{important}

Se precisó que puede existir un curador provisorio mientras se sustancia el proceso de declaración de incapacidad, hasta que el Cuerpo Médico Forense produzca los informes correspondientes.

\begin{quote}
\textit{Docente:} Es el único caso que tiene que subir en consulta a la cámara, aunque no se haya apelado, para confirmar la decisión de la jueza o del juez que lo declaró insano.
\end{quote}

También se discutió el supuesto en que el curador propuesto tiene intereses contrapuestos con la persona incapaz —por ejemplo, si el curador es a la vez heredero en una sucesión en la que también participa la persona bajo curatela—; en tal caso corresponde designar a otro representante, precisamente por existir un conflicto de intereses.

\section{Fuero de atracción y procesos universales}

En los procesos universales —sucesiones, quiebras y concursos— todos los procesos vinculados al patrimonio de la persona física o jurídica involucrada deben ser atraídos y tramitados ante el mismo juez, con el objetivo de concentrar el tratamiento del patrimonio en un único proceso.

\begin{quote}
\textit{Docente:} El fuero de atracción es operativo en los procesos universales. Cuando hay lo mismo, se juntan mismos patrimonios, o sea, el patrimonio de esa persona física o jurídica tiene que ser atraído por el mismo juez o jueza para tratar procesos universales únicos: sucesiones, quiebras, concursos.
\end{quote}

\begin{important}
Conforme fue expuesto en clase, cuando el locatario utiliza el inmueble exclusivamente para vivienda propia o de su familia, el contrato de locación queda ajeno al concurso o quiebra: no opera el fuero de atracción. En cambio, cuando el contrato de locación tiene un destino comercial, sí resulta atraído por el fuero del concurso o la quiebra.

% TODO: contenido incompleto o ambiguo en las notas originales
Nota de cátedra: el docente citó la norma en clase como ``artículo 157 de la ley de concursos y quiebras'' (Ley 24.522). Dado que el texto exacto no pudo verificarse con la misma certeza que el resto de las normas citadas, se recomienda cotejar la referencia puntual con el texto vigente antes de invocarla en un escrito judicial.
\end{important}

\begin{quote}
\textit{Docente:} Si te viene un contrato con destino comercial, acordate que está el fuero de atracción y te vas al juez del concurso y quiebra. Si te viene para desalojar un contrato destino vivienda, ya sabés que queda afuera. Listo, más fácil.
\end{quote}

Se aclaró, finalmente, que esta distinción no siempre es advertida de oficio por el tribunal, y que en la práctica depende en buena medida de que la contraparte la plantee expresamente.

\chapter{Normativa de fondo y causales de desalojo}

\section{El contrato como ley entre las partes}

\begin{example}{La carta documento con norma derogada}
El docente relató que, al recibir un caso, lo primero que solicitó a la persona que consultaba fue el contrato de locación, remarcando la importancia de no confiar en el relato oral del cliente sobre la fecha o el contenido del contrato, sino de exigir siempre el documento original o una copia fehaciente. Al revisar el contrato —celebrado en el año 2025, ya bajo la vigencia plena del DNU 70/2023—, advirtió que la carta documento invocaba un artículo del Código Civil y Comercial que había sido derogado precisamente por ese decreto. Además, la cláusula contractual invocada contenía una renuncia de las partes mayores y capaces a valerse de una norma vinculada a los desperfectos o problemas que afectan el uso y goce del inmueble.
\end{example}

% TODO: contenido incompleto o ambiguo en las notas originales
\begin{important}
En el audio de la clase, el docente identifica esta norma como ``artículo 1223 del Código Civil'' al referirse a la posibilidad de renunciar a un derecho vinculado con desperfectos que impiden el uso y goce del inmueble. Verificado el texto vigente del Código Civil y Comercial, el artículo 1223 CCyC regula en realidad el instituto del ``Desalojo'' en sentido estricto (la restitución de la tenencia del inmueble al extinguirse la locación y el plazo mínimo de ejecución de la sentencia respectiva), y no la cuestión de los desperfectos que afectan el uso y goce de la cosa locada. Esa segunda cuestión se encuentra regulada, dentro del régimen de obligaciones del locador, en las normas sobre conservación de la cosa con aptitud para el uso convenido y garantía por defectos (arts.\ 1200 a 1204 CCyC, en particular el art.\ 1201 CCyC sobre frustración del uso o goce de la cosa). Se deja esta precisión como corrección editorial, sin alterar la referencia oral del docente, a fin de que quien utilice este material verifique el número exacto de artículo antes de invocarlo en un escrito judicial.
\end{important}

Más allá de la imprecisión en el número puntual del artículo, el principio jurídico que el docente quiso transmitir con este caso es central:

\begin{important}
\textbf{Art.\ 958, Código Civil y Comercial de la Nación — Libertad de contratación.} El docente citó este artículo como fundamento del principio de que ``el contrato es ley para las partes'', en referencia a la norma que consagra la libertad de las partes para celebrar un contrato y determinar su contenido, dentro de los límites impuestos por la ley, el orden público, la moral y las buenas costumbres.
\end{important}

\begin{quote}
\textit{Docente:} ¿Qué les dije yo del 958? Que el contrato es ley para las partes. Tengo que ver el contrato, qué firmaron. Hay cláusulas, no es que todo el mundo firma lo mismo; tengo que leer el contrato, lamentablemente es así. Y si no lo leo, hago burradas como la que contesté en la carta documento.
\end{quote}

El docente completó la reflexión señalando que, si una cláusula contractual implica la renuncia a un derecho, es preciso determinar si esa cláusula es válida: existen artículos del Código Civil y Comercial que son de orden público y, por lo tanto, irrenunciables, más allá de lo que las partes hayan pactado. Recae en quien recibe y analiza el contrato la responsabilidad de detectar esas cláusulas potencialmente nulas.

\section{Intimación de pago previa al desalojo por falta de pago}

\begin{important}
\textbf{Art.\ 1222, Código Civil y Comercial de la Nación — Intimación de pago.} ``Si el destino es habitacional, previamente a la demanda de desalojo por falta de pago de alquileres, el locador debe intimar fehacientemente al locatario el pago de la cantidad debida, otorgando para ello un plazo que nunca debe ser inferior a diez días corridos contados a partir de la recepción de la intimación, consignando el lugar de pago. La notificación remitida al domicilio denunciado en el contrato por el locatario se tiene por válida, aun si éste se negara a recibirla o no pudiese perfeccionarse por motivos imputables al mismo.''

Nota de cátedra: el docente advirtió que existe un proyecto legislativo orientado a acotar o directamente eliminar esta exigencia de intimación previa, con el propósito de habilitar más rápidamente la vía judicial del desalojo por falta de pago.
\end{important}

Se precisó que la obligación de intimar fehacientemente rige específicamente para la causal de falta de pago, y que la exigencia de intimación previa no está prevista de igual manera para otras causales de desalojo.

\begin{quote}
\textit{Docente:} La obligación es para falta de pago. La intimación no está prevista para otras causales.
\end{quote}

\section{Canon locativo, expensas y orden público en la falta de pago}

\begin{definition}{Orden público en la falta de pago}
El Código Civil y Comercial exige, como mínimo, dos períodos de canon locativo adeudados para habilitar la vía del desalojo por falta de pago. Se trata de una norma de orden público: el contrato puede pactar un plazo mayor de tolerancia (por ejemplo, tres meses), pero nunca uno menor a los dos períodos legalmente exigidos.
\end{definition}

\begin{quote}
\textit{Docente:} Podría decir el contrato tres, pero no uno. Esto es mejorar, pero no perjudicar, para no violar el orden público.
\end{quote}

El docente advirtió que se trata de una defensa que debe ser invocada por la parte interesada: si el locatario (o su abogado) no plantea que el contrato viola el orden público al exigir un solo mes de mora, el juez no lo declarará de oficio en todos los casos.

\begin{quote}
\textit{Docente:} Si el otro lado no te lo alega, jorobate.
\end{quote}

\begin{example}{Imputación de pagos: expensas y canon locativo}
El docente explicó una recomendación habitual a sus clientes locadores: si el locatario adeuda expensas y paga puntualmente el canon del mes de agosto, corresponde indicarle al cliente que impute ese pago primero a la deuda de expensas (dejando constancia expresa del concepto imputado), de modo que, cuando llegue septiembre sin que se abone el alquiler correspondiente, el locador quede habilitado a demandar por los dos períodos adeudados (agosto y septiembre), ya vencido el plazo contractual de pago (por ejemplo, del día 1 al día 10 de cada mes).
\end{example}

Se aclaró, además, que la obligación de pagar las expensas es una obligación \textit{propter rem} —que sigue a la cosa—, motivo por el cual, si el consorcio de propietarios decide reclamar expensas impagas, dirige su ejecución contra el titular de dominio del inmueble y no contra el ocupante o locatario, más allá de que en la relación interna entre locador y locatario este último se haya comprometido contractualmente a afrontarlas.

El docente cerró esta sección con una reflexión sobre la fijación libre del canon locativo, tras la derogación de la Ley de Alquileres por el DNU 70/2023: el precio del alquiler queda sujeto, en principio, al valor de mercado y a la libre negociación entre locador y locatario, pudiendo incluir o no las expensas dentro del monto total pactado.

\section{Recusación sin causa en el juicio de desalojo}

\begin{note}
\textbf{Art.\ 17, Código Procesal Civil y Comercial de la Nación — Recusación con causa.} El docente hizo referencia a este artículo como el que regula la recusación con expresión de causa —cuando existe un motivo concreto, como amistad o enemistad con alguna de las partes, que justifica apartar al juez del conocimiento de la causa—, distinguiéndolo del instituto que se desarrolla a continuación.
\end{note}

\begin{important}
\textbf{Art.\ 14, Código Procesal Civil y Comercial de la Nación — Recusación sin expresión de causa.} ``Los jueces de primera instancia podrán ser recusados sin expresión de causa. El actor podrá ejercer esta facultad al entablar la demanda o en su primera presentación; el demandado, en su primera presentación, antes o al tiempo de contestarla, o de oponer excepciones en el juicio ejecutivo, o de comparecer a la audiencia señalada como primer acto procesal (\ldots). No procede la recusación sin expresión de causa en el proceso sumarísimo ni en las tercerías, en el juicio de desalojo y en los procesos de ejecución.''

Nota de cátedra: la última oración fue incorporada por la Ley 25.488 (B.O.\ 22/11/2001), que extendió la prohibición de recusar sin causa —antes limitada al juicio ejecutivo— al proceso sumarísimo, a las tercerías, al juicio de desalojo y a los procesos de ejecución.
\end{important}

El docente explicó que esta prohibición se aplica de manera directa y expresa al juicio de desalojo, por texto literal del propio artículo 14. Se presentó, además, un debate jurisprudencial conexo: si la prohibición de recusar sin causa en los ``procesos de ejecución'' alcanza también a la ejecución de alquileres.

\begin{quote}
\textit{Estudiante:} ¿Por qué dice el artículo 14 esto de que no opera del\ldots del 2003?

\textit{Docente:} El artículo 14 te dice esto, pero también te agregué que el artículo 14 te dice que no procede en los procesos ejecutivos. Y como dijo tu compañero, ¿qué hacemos si debe alquileres o debe otros rubros? Hay que empezar otro juicio que se llama ejecución de alquileres o cobro de alquileres, que verán ustedes en la última unidad.
\end{quote}

\begin{note}
El docente explicó que, si bien el texto literal del artículo 14 no menciona expresamente la ``ejecución de alquileres'', tanto la Cámara Nacional de Apelaciones en lo Comercial en pleno (fallo plenario del 27 de mayo de 2003, ``Recusación sin causa en juicios ejecutivos s/convocatoria a plenario'') como la Cámara Nacional de Apelaciones en lo Civil, en criterio análogo, resolvieron que la prohibición de recusar sin expresión de causa contenida en el artículo 14 del Código Procesal alcanza también a los juicios ejecutivos, categoría dentro de la cual se ha considerado comprendida la ejecución de alquileres.

Conclusión práctica: pese a que los abogados de la contraparte a veces igualmente intentan la recusación sin causa en estos procesos, dicho planteo debe ser rechazado por el juzgado en virtud de la doctrina fijada en los plenarios citados.
\end{note}

\chapter{Debido proceso y herramientas prácticas}

\section{El debido proceso legal: faceta adjetiva y sustantiva}

\begin{definition}{Faceta adjetiva y faceta sustantiva del debido proceso}
La faceta adjetiva (o formal) exige que se cumplan todos los recaudos procesales previstos por la ley: el traslado de la demanda, el plazo para ejercer la debida defensa, y la notificación correcta al domicilio en el que efectivamente se llevará a cabo el desalojo —incluyendo la precisión sobre si se trata del piso de arriba, de abajo, de atrás o de al lado—. Si estos recaudos no se cumplen, la persona nunca pudo ejercer su defensa y no puede considerarse válidamente notificada.

La faceta sustantiva (o material) exige que, a lo largo de todo el proceso, se respeten no solo la Constitución Nacional —que el docente calificó como ``operativa''— sino también los tratados internacionales de derechos humanos a los que la Argentina se ha adherido y continúa adhiriendo, lo cual implica que el catálogo de garantías a tener en cuenta se amplía progresivamente con cada nuevo tratado incorporado.
\end{definition}

El docente explicó que en su cátedra se diferencia de manera especial entre principios y sistemas: los principios provienen del poder constituyente (la Constitución) y del poder convencional (los tratados internacionales), mientras que los sistemas —como el propio Código Procesal— son creación del legislador ordinario y deben, necesariamente, adecuarse a los principios.

\begin{quote}
\textit{Docente:} La Constitución Nacional es un compendio de principios. El tratado de Belém do Pará es un principio contra la vulnerabilidad de la mujer. Los sistemas tienen que adecuarse a los principios que rigen en Argentina. Singapur tendrá otros principios, seguro que sí; Brasil tendrá otros principios. Los sistemas bajan a la realidad del proceso lo que los principios rigen en este país, Argentina.
\end{quote}

Se mencionó, dentro de este eje, la relevancia de revisar si el contrato contiene una cláusula de prórroga de competencia (pacto de foro prorrogado) o una cláusula compromisoria que derive la conflictiva hacia el arbitraje, así como el instituto del fuero de atracción ya desarrollado en el capítulo anterior, en tanto ambos son ejemplos de cómo la competencia del juez natural puede verse desplazada.

\section{Métodos alternativos de resolución de conflictos}

El docente dedicó una parte sustancial del cierre de la clase a las técnicas alternativas de resolución de conflictos, presentadas explícitamente como un complemento —y no como un reemplazo— de la vía judicial:

\begin{quote}
\textit{Docente:} Acuérdense que yo soy hincha de que está la vía procesal judicial, pero no saquemos de nuestro nicho mental las otras vías; si no, hagámoslas en paralelo, porque el proceso judicial es fatal.
\end{quote}

Entre las técnicas desarrolladas se destacaron las siguientes:

\begin{itemize}
\item \textbf{Escucha activa:} el docente distinguió entre simplemente oír y escuchar activamente, y explicó su preferencia personal por la comunicación escrita por sobre los mensajes de audio, precisamente para poder documentar con precisión lo indicado por cada cliente y evitar reproches posteriores sobre instrucciones mal comunicadas.
\item \textbf{Análisis de los intereses en juego:} se recomendó explicar siempre al cliente los costos y los riesgos de iniciar un proceso judicial, incluyendo la posibilidad de que las costas se inviertan y deban afrontarse.
\item \textbf{Uso de parámetros objetivos:} el docente mencionó, como ejemplo, el recurso a una pericial arbitral (distinta de la prueba pericial dentro de un proceso judicial) para dirimir cuestiones técnicas puntuales entre las partes.
\item \textbf{Ampliación de las opciones disponibles} (``agrandar la torta de recursos''): identificar alternativas creativas antes de litigar, en lugar de limitarse a una única solución posible.
\end{itemize}

\begin{example}{La pericial arbitral y el calefón}
El docente relató el caso de un locatario que se negaba a pagar el alquiler argumentando que el calefón del inmueble no funcionaba correctamente. Ante ese planteo, propuso designar un plomero como perito arbitral para que examinara el artefacto. El profesional constató que el desperfecto se debía a un uso incorrecto de la palanca reguladora entre agua fría y agua caliente por parte del propio ocupante, y no a un defecto atribuible al locador. La pericial arbitral permitió, en este caso, despejar objetivamente el reclamo y evitar una discusión judicial prolongada.
\end{example}

El docente desarrolló también la técnica de ``una cosa por otra'' —un intercambio de tipo compensatorio— aplicada a la negociación de la restitución anticipada del inmueble: ofrecer al ocupante una suma de dinero para solventar su mudanza, a cambio de que se retire voluntariamente y sin necesidad de afrontar un proceso judicial que puede extenderse más de un año y medio.

\begin{quote}
\textit{Docente:} Vos le das para la mudanza, dejale que te deba los dos meses de alquiler que te debe, pero que se vaya. Vamos a tener menos de un año y medio de proceso judicial seguro.
\end{quote}

\section{Garantías locativas y su ejecución práctica}

Respecto de la figura del garante o fiador, el docente advirtió sobre un riesgo práctico frecuente: la garantía puede haberse constituido varios años antes de que el conflicto se materialice, período durante el cual el garante puede haber transferido la titularidad de los bienes que originalmente respaldaban la garantía. Recomendó, por ello, solicitar siempre un informe de dominio o de titularidad actualizado del garante antes de aceptar la garantía y, eventualmente, también al momento de iniciar su ejecución.

\begin{quote}
\textit{Docente:} Con el garante podemos tener hermoso en el contrato que hicimos en el 2024, y le pedimos un informe de dominio; resulta que cambió la titularidad o lo vendió, y después le pedimos informes de titularidad, y ese garante no tiene ni la bicicleta que usa.
\end{quote}

\begin{definition}{Tipos de garantía mencionados en clase}
\textbf{Fianza:} garantía personal por excelencia en las locaciones, en la que el fiador se obliga a responder por las deudas del locatario, sin que la obligación siga a un bien determinado (a diferencia de una garantía real).

\textbf{Prohibición personal de disponer:} el docente aclaró que no debe confundirse con una hipoteca —que es una garantía real sobre un inmueble determinado y que sigue a la cosa—, ya que la prohibición personal no tiene ese efecto de persecución sobre el bien.

\textbf{Inhibición general de bienes:} hoy se practica exclusivamente como medida cautelar ordenada por un juez, previo cumplimiento de los presupuestos clásicos de toda cautelar (verosimilitud del derecho y peligro en la demora), y sujeta a contracautela. La llamada ``autoinhibición'' voluntaria de bienes —que la persona consentía en forma autónoma— existió como figura hace más de un siglo, pero fue derogada y no se practica en la actualidad, sin que exista proyecto legislativo alguno para reintroducirla.
\end{definition}

Frente a una consulta sobre la posibilidad de instrumentar algún trámite administrativo adicional para reforzar la garantía sobre el inmueble del garante (impidiendo su venta), el docente fue categórico en cuanto a los límites de esa vía: cualquier intento de calificar como insolvencia fraudulenta la venta de un bien por parte del garante requeriría acreditarlo en sede penal, lo cual suele implicar costos superiores al beneficio económico que se busca proteger.

Como cierre práctico de la clase, el docente compartió su recomendación personal para reforzar las garantías locativas en los contratos que asesora: exigir, además del depósito habitual de garantía (equivalente, por lo general, a un mes de alquiler por año de contrato, o a un mes fijo), una caución real adicional.

\begin{example}{La caución real como incentivo a la negociación}
El docente relató que, en los contratos en los que se pactó que el locatario o su fiador entreguen al locador una caución real de una suma determinada de dinero (mencionó, a título de ejemplo, una suma de \$3.000, en referencia a valores históricos), rara vez se llega a juicio: las partes prefieren sentarse a negociar y resolver el conflicto —incluso cuestiones menores, como el desperfecto del calefón del ejemplo anterior— antes que perder esa suma entregada en garantía. ``Es raro. Viene a sentarse a la mesa a ponerse de acuerdo.''
\end{example}

El docente cerró remarcando que se trata de una práctica de gestión contractual y no de una exigencia legal, cuya eficacia se apoya en el incentivo económico que genera en ambas partes para evitar la vía judicial.

\chapter{Cuadro Cornell y síntesis final}

La siguiente tabla organiza, en formato Cornell, las preguntas y conceptos clave desarrollados a lo largo de la clase junto con su respuesta o explicación correspondiente. La última fila contiene la síntesis integradora de todo el contenido abordado.

\begin{longtable}{p{0.30\textwidth} p{0.60\textwidth}}
\toprule
\textbf{Preguntas / palabras clave} & \textbf{Notas / respuestas} \\
\midrule
\endfirsthead
\toprule
\textbf{Preguntas / palabras clave} & \textbf{Notas / respuestas} \\
\midrule
\endhead
\midrule
\multicolumn{2}{r}{\textit{Continúa en la página siguiente}} \\
\endfoot
\bottomrule
\endlastfoot

\textbf{¿Qué es una diligencia preliminar?} &
Una medida que puede solicitar quien pretende demandar o quien prevé, con fundamento, que será demandado, con el fin de preparar el proceso de conocimiento antes de iniciar la demanda (art.\ 323 CPCCN y ss.). Está sujeta a un plazo de caducidad de 30 días para iniciar el proceso principal. \\

\textbf{¿En qué se diferencia de la prueba anticipada?} &
La diligencia preliminar prepara el proceso antes de demandar y tiene carácter restrictivo y residual. La prueba anticipada es un sistema asegurativo probatorio que evita la pérdida de una prueba (por ejemplo, un testigo próximo a fallecer) y constituye un instituto autónomo, no subordinado a la diligencia preliminar. \\

\textbf{¿Qué dice el art.\ 323 inc.\ 1º CPCCN?} &
Habilita a pedir que la persona contra quien se proponga dirigir la demanda preste declaración jurada, por escrito, sobre algún hecho relativo a su personalidad, sin cuya comprobación no pueda entrarse en juicio: herramienta clave para identificar al ocupante de un inmueble antes de demandar. \\

\textbf{¿Cuál es el plazo de caducidad de la diligencia preliminar?} &
30 días corridos desde su sustanciación para iniciar el proceso principal. Si no se inicia en ese plazo, la diligencia pierde sus efectos. \\

\textbf{¿Cómo se diferencian discovery/disclosure del sistema argentino?} &
En Estados Unidos (discovery) y en Gran Bretaña (disclosure) existe una etapa de preparación del proceso con un fuerte componente extrajudicial; en la Argentina, las diligencias preliminares son netamente judiciales. \\

\textbf{¿Qué procesos de conocimiento existen en el orden nacional?} &
Ordinario y sumarísimo. La mayoría de los desalojos tramita por la vía sumarísima. \\

\textbf{¿Qué procesos de conocimiento existen en la provincia de Buenos Aires?} &
Ordinario, sumario y sumarísimo: los tres subsisten, a diferencia del orden nacional, donde el proceso sumario fue eliminado por la Ley 25.488. \\

\textbf{¿Qué diferencia hay entre legitimación sustancial y legitimación procesal?} &
La legitimación sustancial es ser el titular de la relación jurídica de fondo (equivalente a la capacidad de derecho). La legitimación procesal es la aptitud para actuar por sí mismo en el proceso (equivalente a la capacidad de ejercicio, art.\ 24 CCyC). Puede tenerse una sin la otra, como en el caso de una persona declarada incapaz por sentencia judicial. \\

\textbf{¿Cuándo opera el fuero de atracción respecto de un desalojo?} &
Cuando el contrato de locación tiene destino comercial y existe un concurso o quiebra del locatario. Si el destino es habitacional, el contrato queda ajeno al concurso o quiebra y no opera el fuero de atracción. \\

\textbf{¿Qué significa que ``el contrato es ley para las partes''?} &
Principio consagrado en el art.\ 958 CCyC: obliga a leer cada cláusula del contrato con atención, ya que las partes pueden haber pactado condiciones específicas, y a verificar si alguna cláusula pretende renunciar a un derecho de orden público, lo que la tornaría inválida. \\

\textbf{¿Qué exige el art.\ 1222 CCyC antes de demandar por falta de pago en un inmueble habitacional?} &
Una intimación fehaciente previa al locatario, otorgando un plazo no inferior a diez días corridos desde la recepción de la intimación, con indicación del lugar de pago. \\

\textbf{¿Cuántos períodos de falta de pago se requieren para habilitar el desalojo?} &
Como mínimo, dos períodos consecutivos de canon locativo adeudados, por tratarse de una norma de orden público que el contrato puede ampliar pero nunca reducir. \\

\textbf{¿Procede la recusación sin causa en el juicio de desalojo?} &
No. El art.\ 14 CPCCN (según Ley 25.488) excluye expresamente el juicio de desalojo, el proceso sumarísimo, las tercerías y los procesos de ejecución de la posibilidad de recusar sin expresión de causa. \\

\textbf{¿La ejecución de alquileres está alcanzada por esa prohibición?} &
Sí, según la doctrina fijada por sendos plenarios de la Cámara Nacional de Apelaciones en lo Comercial (27/05/2003) y de la Cámara Nacional de Apelaciones en lo Civil, que interpretaron que la ejecución de alquileres queda comprendida dentro de los ``procesos de ejecución'' mencionados por el art.\ 14 CPCCN. \\

\textbf{¿Qué dos facetas comprende el debido proceso legal?} &
La faceta adjetiva o formal (notificación correcta, traslado, plazo de defensa) y la faceta sustantiva o material (respeto de la Constitución y de los tratados internacionales de derechos humanos, en constante ampliación). \\

\textbf{¿Qué diferencia hay entre principios y sistemas, según la cátedra?} &
Los principios provienen del poder constituyente y convencional (Constitución y tratados internacionales); los sistemas son creación del legislador (por ejemplo, el Código Procesal) y deben adecuarse siempre a los principios. \\

\textbf{¿Qué técnicas de resolución alternativa de conflictos se desarrollaron en clase?} &
Escucha activa, análisis de intereses y de costos/riesgos, uso de parámetros objetivos (por ejemplo, pericial arbitral) y ampliación de las opciones disponibles antes de litigar (``agrandar la torta de recursos''). \\

\textbf{¿Qué garantías locativas se recomendaron en clase?} &
Depósito de garantía habitual, fianza personal, y —como recomendación práctica del docente— una caución real adicional, que en la experiencia relatada desalienta eficazmente los conflictos judiciales al generar un incentivo económico a la negociación. \\

\midrule
\multicolumn{2}{p{0.90\textwidth}}{
\textbf{Resumen:} El proceso de desalojo, lejos de ser un trámite meramente mecánico, exige del abogado una preparación estratégica previa: identificar con precisión al ocupante del inmueble mediante diligencias preliminares (art.\ 323 CPCCN y ss.), sin confundir este instituto con la prueba anticipada, y siempre dentro del plazo de caducidad de treinta días. Requiere, además, conocer con exactitud qué proceso de conocimiento corresponde según la jurisdicción (ordinario o sumarísimo en el orden nacional; ordinario, sumario o sumarísimo en la provincia de Buenos Aires) y verificar la legitimación —tanto sustancial como procesal— de todas las partes involucradas, incluyendo los supuestos de incapacidad de ejercicio regulados en el art.\ 24 CCyC. El operador jurídico debe, asimismo, evaluar si el destino del contrato locativo (habitacional o comercial) desplaza la competencia hacia el fuero de atracción de un concurso o quiebra, y tener presente que ni el desalojo ni la ejecución de alquileres admiten la recusación sin expresión de causa del juez interviniente (art.\ 14 CPCCN). En el plano sustantivo, el eje que atraviesa toda la clase es la primacía del contrato como ley entre las partes (art.\ 958 CCyC) y la necesidad de leerlo íntegramente antes de invocar cualquier norma, verificando siempre su vigencia frente a reformas sucesivas como el DNU 70/2023, que modificó buena parte del régimen locativo del Código Civil y Comercial. Finalmente, la clase deja una enseñanza que excede lo estrictamente procesal: el debido proceso legal —en sus facetas adjetiva y sustantiva— y los métodos alternativos de resolución de conflictos (escucha activa, análisis de intereses, pericial arbitral, garantías reales como incentivo a la negociación) no son un complemento accesorio del litigio, sino herramientas centrales para ejercer una abogacía eficaz, económica y más humana en la gestión de los conflictos locativos.
} \\

\end{longtable}

\end{document}
